\abstract{%
Intrinsic-dimension estimates and active-unit diagnostics address different
aspects of variational-autoencoder bottleneck selection. We evaluate
search policies on four image datasets, five training seeds and five
regularisation weights, using a validation quality target and common
recorded candidate costs. A start-width ablation holds reliability checks,
fallback and reference-bank charges fixed. ID is accepted only for the
25 MNIST conditions: ID-derived and gated midpoint starts select the
same widths, requesting 8.84 and 5.84 candidates on average, respectively.
The remaining 75 conditions use the same ascending fallback. Across
10,000 replay evaluations from 100 coefficient--threshold combinations
and 100 dataset--weight--seed conditions, 155 local outputs miss the
smallest qualifying width. Three-seed pruning experiments per dataset
distinguish reconstruction by the selected native model from dimension
transfer to an ordinary variational autoencoder. In a local relevance-prior
variant, masking nonselected axes increases dSprites validation error
by 126.3\%; 30 decoder-adaptation epochs reduce this to 37.1\% above
the full model. All four evaluated conditions, including the full model,
miss that dataset's quality target. Active-unit prediction does not
establish additional benefit on its near-ceiling task. The findings
support a controlled empirical evaluation of selection rules and
diagnostics, without a general efficiency or pruning guarantee.}
\keyword{variational autoencoder; bottleneck dimension; intrinsic dimension;
active units; model selection; reconstruction quality; reproducibility}
